\section{Parallel Patterns}

\subsection{Divide and Conquer}

\begin{definition}
    \textbf{Divide and Conquer } refers to the strategy of splitting a problem into smaller subproblems, solving them recursively, and combining the results.
\end{definition}

This concept can be used in many parallel algorithms. Until now we have used manual forking .start() and joining .join() of Threads. We can continue using manual fork/join with DC, but this can lead to several Problems.

\begin{itemize}
    \item Poor reusability of code
    \item We need to commit to a number of threads a priori and can not dynamically adjust.
    \item Subproblems are not automatically assigned to a thread.
\end{itemize}

\begin{remark}
    There are some fixes that can improve manual fork/join such as \textbf{sequential-cutoffs} or running one subproblem on the thread itself (\textbf{invoke/compute}), instad of creating two new threads.
\end{remark}

\subsubsection{Executor Service}

The Executor Service framework provides the way to schedule and run tasks on a thread-pool automatically.

\begin{definition}
    \textbf{Runnable}: A task without a return value. Passed to execute() (fire-and-forget) or submit() (returns an empty future).
\end{definition}

\begin{definition}
    \textbf{Callable}: A task that returns a result of given type T and can throw checked exceptions. Passed to submit(Callable) and returns a Future of type T.
\end{definition}

\begin{remark}
    There are diffrent \textbf{Pool Types}: Executors.newFixedThreadPool(n) keeps $n$ worker threads alive. newCachedThreadPool() scales up and down dynamically depending on load.
\end{remark}

\begin{codeexample}
ExecutorService executor = Executors.newFixedThreadPool(4);

class MyRunnable implements Runnable {
    public void run() {...}
}

class MyCallable implements Callable<Integer> {
    public Integer call() throws Exception {...}
}   

executor.execute(new MyRunnable());

Future<Integer> future = executor.submit(new MyCallable());
int res = future.get(); // Blocks until result is returned

executor.shutdown();
\end{codeexample}

Pros and Cons of Executor service:
\begin{itemize}
    \item Not suited for: Recursive subproblems that depend on each other and have a deep fork-tree. As the service will eventually run out of threads.
    \item Suited for: flat structures that can run independently in parallel.
\end{itemize}

\subsubsection{Fork-Join}

The fork-join framework is built for divide and conquer parallelism. It manages a pool of threads using a \textbf{work-stealing strategy}: each thread maintains a local task-queue, and when finished with its own tasks, steals work from others to minimize idle time.

\begin{definition}
    \textbf{RecursiveTask}: A task that returns a result of type given Type T (equivalent to Callable).
\end{definition}

\begin{definition}
    \textbf{RecursiveAction}: A task that does not return a result (equivalent to Runnable).
\end{definition}

\textbf{Key Methods of the Fork-Join Framework:}
\begin{itemize}
    \item \textbf{fork()}: Asynchronously pushes the task to the local thread's queue.
    \item \textbf{join()}: Blocks until the result is ready. While waiting, the thread does not truly idle but uses work-stealing to execute other tasks.
    \item \textbf{compute()}: The main logic method that must be overridden.
\end{itemize}

\begin{codeexample}
class MyTask extends RecursiveTask<Integer> {
    int start, end;
    MyTask(int start, int end) { this.start = start; this.end = end; }
    
    protected Integer compute() {
        if (end - start <= SEQUENTIAL_CUTOFF) { 
            return baseCase(); 
        } else {
            int mid = (start + end) / 2;
            MyTask left = new MyTask(start, mid);
            MyTask right = new MyTask(mid, end);
            
            left.fork(); // Async execute left half
            int rightRes = right.compute(); // Optimize: caller thread computes right half
            int leftRes = left.join();      // Wait for left half
            
            return leftRes + rightRes;
        } 
    }
}
\end{codeexample}

\begin{important}
    Make sure that every task has enough work to do (aprox. 100-10'000 ops.), else the overhead will be larger than the parallel speedup gain.
\end{important}

\subsubsection{Task graphs}

The execution of parallel programs using Fork-Join can be represented as a Task-Graph (DAG).

\begin{definition}
    We introduce a new size $T_{\infty}$ which is the execution time assuming unlimited hardware. This is called the \textbf{span} and can be visualized as the sum of the execution times of the longest dependent path in our DAG.
\end{definition}

\begin{remark}
    The maximum speedup is limited by $\frac{T_1}{T_{\infty}} (Amdahls Law)$.
\end{remark}

It is the programmers responsibility to find the optimal decomposition/alogrithm for the problem to ensure a
minimal span (a balanced tree). \textbf{Fork-join will guarantee optimal scheduling of the tasks with minimal idletime.}


\subsection{Parallel Patterns}

\begin{definition}
    A \textbf{reduction} is a parallel operation that produces a single answer from a collection via an
    associative operator using for example DC. (The order of operations does not matter).
\end{definition}

\begin{example}
    Finding the Max, Min, Sum or Product over all entries of an array.  
\end{example}

\begin{definition}
    A \textbf{map} applies a function to each elements of a collection producing an output of the same size.
\end{definition}

\begin{example}
    Vector addition can be done using a map pattern. We split the vectors into dim(v)/2 recursively
    until we have dim(v)=1 then sum the vectors and write the result directly into the result vector.
\end{example}

\begin{remark}
    Maps and Reductions have (when parallelized) an optimal span of $O(\log n)$.
\end{remark}

\begin{definition}
    A \textbf{stencil} is an extension of the map pattern that applies a function to each element of a collection producing an output of the same size. The function uses the element itself and its neighbours.
\end{definition}

\begin{example}
    A smoothing kernel (ex. 3x3) that calculates the average value is sled over an image to smooth/blur the image.
\end{example}

\subsection{Parallel Algorithms}

\subsubsection{Prefix-Sum}

\begin{remark}
    Prefix-Sum (calculating input[1] + input[2] + ... + input[n]) seams to be unparallisable. But infact it is!
\end{remark}

\begin{algorithm}
upTask(); //divide problem and build tree with sums of subproblems
downTask(); //fill in the correction values (left_sum + parent_correction)
\end{algorithm}

\begin{theorem}
    The span of the prefix sum algorithm is $O(\log n)$. This leads to a parallelism of $\frac{n}{\log n}$.
\end{theorem}

\subsubsection{Pack}

Another application of this problem is the \textbf{Pack} algorithm which removes all elements from an array that satisfy a certain condition and packs them into a seperate output.

\begin{algorithm}
applyMap(); //applies a map to the input
//results in a binary vector [1, 0, 1, 1, 0, ...]

parallelPrefixSum(); //prefix sum of the binary vector
//results in [1, 1, 2, 3, 3, ...]

applyPackMap(); //applies map to the result of the prefix sum
//results in [0, 2, 3, ...] (indices of output elements)
\end{algorithm}

\subsubsection{Quicksort}

We can reduce the limiting factor of the problem division from a span of $O(n)$ to a span of $O(\log n)$
using the pack-pattern we have seen before.

\begin{algorithm}
applyMaps(); //apply maps for greater and smaller than pivot
parallelPrefixSum(); //num. of smaler/greater elems. come bofore
applyMap(); //index = numLessThan || numTotalLessThan + numGreaterThan
\end{algorithm}

We can do this by using two pack-patterns to partition the data. This leads to a partitioning of $O(\log n)$
levels times $O(\log n)$ for prefixSum hence $O(\log^2 n)$ total span. This means a parallelism of 
$\frac{n \log n}{\log^2 n} = \frac{n}{\log n}$.

        